\begin{table}[!htbp]
\centering
\footnotesize
\makebox[\textwidth][c]{%
\begin{threeparttable}
\caption{Number of states with CES respondents and mean respondents per state, by office and year.}
\label{tab:states_respondents_by_office}
\begin{tabular}{lcccccccc}
\toprule
 & \multicolumn{8}{c}{Office} \\
\cmidrule(lr){2-9}
Year & President & US Senate & US House & Governor & Atty. gen. & Sec. of state & State sen. & State rep. \\
\midrule
\multicolumn{9}{c}{Number of states with CES respondents} \\
\addlinespace[0.2em]
2006 & n/a & -- & -- & -- & -- & -- & 42 & 44 \\
2008 & 51 & 33 & 48 & -- & -- & -- & 39 & 41 \\
2010 & n/a & 36 & 51 & 37 & -- & -- & 41 & 44 \\
2012 & 51 & 33 & 51 & 11 & -- & -- & 39 & 43 \\
2014 & n/a & 34 & 51 & 35 & -- & -- & 41 & 45 \\
2016 & 51 & 34 & 51 & 12 & -- & -- & 40 & 43 \\
2018 & n/a & 33 & 51 & 36 & 29 & 25 & 40 & 46 \\
2020 & 51 & 34 & 51 & 11 & 9 & 7 & 41 & 44 \\
2022 & n/a & 34 & 51 & 36 & 28 & 23 & 41 & 46 \\
\addlinespace[0.6em]
\multicolumn{9}{c}{Mean number of CES respondents per state} \\
\addlinespace[0.2em]
2006 & n/a & -- & -- & -- & -- & -- & 405 & 410 \\
2008 & 469 & 351 & 422 & -- & -- & -- & 408 & 453 \\
2010 & n/a & 784 & 708 & 807 & -- & -- & 702 & 732 \\
2012 & 769 & 856 & 728 & 492 & -- & -- & 653 & 706 \\
2014 & n/a & 507 & 636 & 773 & -- & -- & 626 & 651 \\
2016 & 878 & 932 & 816 & 547 & -- & -- & 734 & 811 \\
2018 & n/a & 940 & 793 & 900 & 934 & 747 & 809 & 826 \\
2020 & 888 & 699 & 837 & 557 & 974 & 701 & 814 & 891 \\
2022 & n/a & 861 & 762 & 875 & 1058 & 821 & 784 & 800 \\
\bottomrule
\end{tabular}
\begin{tablenotes}
\footnotesize
\item n/a Office not on ballot that year.
\item -- CES did not measure vote choice for that office that year.
\item \textit{Note:} Washington D.C. is counted as the 51st state. The US House counts from 2010 on include the election for the D.C. nonvoting delegate. For US House, respondent counts are summed across congressional districts within each state before averaging.
\end{tablenotes}
\end{threeparttable}%
}
\end{table}